\documentclass{article}
\usepackage{tikz}
\usetikzlibrary{angles, quotes}
\usepackage{amssymb}
\usepackage{pgfplots}
\pgfplotsset{width=10cm,compat=1.9}
\usepgfplotslibrary{statistics}
\usepackage[left=2.54cm,right=2.54cm,top=2.54cm,bottom=2.54cm]{geometry}
\usepackage{tasks} %Paquete necesario para la producción de items horizontales para algunos ejercicios de tarea empleando el comando /begin{multicols}}{#}.../end{multicols} para ayudarnos a reducir las páginas
\usepackage{pdflscape} %comando que permite cambiar a orientación horizontal las hojas%
\usepackage{soul} %Paquete para permitir la compilación de acentos usando el comando {\hl{}}}\\
\sethlcolor{yellow(munsell)} %comando para definir el color para el subrayado usando el comando \hl
\usepackage[utf8]{inputenc}
\usepackage[spanish]{babel}
\usepackage{icomma}
\usepackage{ragged2e}
\usepackage{siunitx}
\usepackage{url}
\usepackage[colorlinks=true, urlcolor=blue,  linkcolor=black, citecolor=green]{hyperref}
\usepackage{pdfpages}
\usepackage{blindtext} %Paquete que genera texto ficticio (dummy text) con el comando: \blindtext
\usepackage{cancel} %in the preamble gives you four different modes of striking through: \cancel{text to cancel} draws a diagonal line (slash) through its argument, \bcancel{text to cancel} uses the negative slope (a backslash), \xcancel{text to cancel} draws an X (actually \cancel plus \bcancel), \cancelto{〈value〉}{〈expression〉} draws a diagonal arrow through the 〈expression〉pointing to the 〈value〉 (math-mode only)
\usepackage{csquotes}
\usepackage{afterpage}
\usepackage{parskip} 
\usepackage{float}
\usepackage{enumitem}
\usepackage{multicol}%Paquete que permite la creación aislada de columnas de texto. De esta forma se reduce la cantidad de páginas en nuestro documento
\usepackage{lipsum} %Paquete que genera texto de relleno al igual que \usepackage{blindtext}, se usa con el comando \lipsum[#-#] los #(hashtags) delimitan el número de páginas que deseamos ocupar del paquete lipsum para usarlos como relleno. 

% Definir el comando personalizado para la nota
\newcommand{\mynote}[1]{%
\begin{tikzpicture}[baseline=-0.75ex]
\node[draw, circle, fill=Apple Green!60, inner sep=2pt] (note) {\textbf{Nota:}};
\end{tikzpicture}%
\ \textit{#1}%
}

\newenvironment{Figura}
{\par\medskip\noindent\minipage{\linewidth}}
{\endminipage\par\medskip}
\usepackage{caption}
\usepackage[
backend=biber,
sorting=none,
url=true
]{biblatex} %bibliografía
\addbibresource{biblio.bib}
% Establecer el espacio entre las entradas de la bibliografía
\setlength{\bibitemsep}{1\baselineskip} % Puedes ajustar el valor según tus preferencias
\usepackage{amsmath,amsthm,amssymb,amsfonts}
\usepackage{pifont} %Permite la compilación del comando \xmark: es el símbolo de la tachita
\usepackage{mathtools} %permite la compilación de símbolos para matrices
\usepackage{empheq} %Paquete que se relaciona con \usepackage[most]{tcolorbox} para la creación de cuadros/cajas de colores para delimitar los resultados matemáticos o líneas de texto usando la línea de comando: \begin{empheq}[box={\mymath[colback=orange(webcolor)!70,drop lifted shadow]}]{equation*} \end{empheq}
\usepackage[most]{tcolorbox} %Paquete que se relaciona con \usepackage{empheq} para la creación de cuadros/cajas de colores para delimitar los resultados matemáticos o líneas de texto
\renewcommand{\qedsymbol}{$\blacksquare$}
\usetikzlibrary{positioning,decorations.pathreplacing} %Paquete que permite la compilación de llaves para cuadros sinópticos (brace diagram) 
\usepackage{schemata} %The schemata package is designed just for these brace diagrams
\usepackage{booktabs}

\newcommand\AB[2]{\schema{\schemabox{#1}}{\schemabox{#2}}} %comando o paquete necesario para crear cuadros sinópticos

\newtcbox{\mymath}[1][]{%
nobeforeafter, math upper, tcbox raise base,
enhanced, colframe=blue!30!black,
colback=blue!30, boxrule=1pt,
#1} %Comando importante para encasillar los resultados matemáticos en cajas de diferentes colores, se relaciona con los paquetes \usepackage{empheq} y \usepackage[most]{tcolorbox} 

\newenvironment{sysmatrix}[1]
{\left(\begin{array}{@{}#1@{}}}
{\end{array}\right)}
\newcommand{\ro}[1]{%
\xrightarrow{\mathmakebox[\rowidth]{#1}}%
}
\newlength{\rowidth}% row operation width
\AtBeginDocument{\setlength{\rowidth}{3em}}  %Comando importante para laproducción de lineas de operaciones entre matrices, método de Gauss o Gauss Jordan se relaciona con \newenvironment{sysmatrix}

%formato para cambiar el horario a español
\usepackage[spanish]{datetime2}
\DTMsetdatestyle{spanish}

\renewcommand{\today}{\DTMdisplaydate{\the\year}{\the\month}{\the\day }{-1}}
%%%%%%%%%%%%%%%%%%%%%%%%%%%%%%%%%%%%%%%%%

\usepackage{datetime}
\newcommand{\mycurrenttime}{\xxivtime}

%%%%%%%%%%%%%%%%%%%%%%%%%%%%%%%%%%%%%%%%%%%%%%%%%%%%%%%%%%%%%%%%%%%%%%%%%%%%%%%%%%%%%%%%%%%%%%%%%%%%%%%%%%%%%%%%%%%%%%%%%%%%%%%%%%%%%%%%%%%%%%%%%


%%%%%%%%%%%%%Caja de color para el título%%%%%%%%%%%%%%%%%%%%%%%%%%%%%%%%%%%%%%%%%%%%%%%%%%%%%%

\definecolor{myframecolor}{RGB}{1, 45, 75} % Prussian Blue
\definecolor{myboxcolor}{RGB}{0, 107, 92} % Apple Green

\newtcolorbox{mybox}{
enhanced,
colback=myboxcolor!7, % Color de fondo del cuadro
colframe=myframecolor, % Color del marco del cuadro
arc=0pt,
boxrule=1pt,
borderline west={2mm}{-10mm}{myframecolor}, % Borde en el lado izquierdo
sharp corners=southwest,
width=\linewidth,
before=\par\vspace{\bigskipamount}, % Espacio antes del cuadro
after=\par\vspace{\bigskipamount} % Espacio después del cuadro
}


\newcommand{\euler}{e} %Comando para producir letra e de euler
\newenvironment{remark}{\par\vfill\footnotesize % Vertical white space above the remark and smaller font size
\begin{list}{}{
	\leftmargin=80pt % Indentation on the left
	\rightmargin=60pt}\item\ignorespaces % Indentation on the right
\makebox[-2.5pt]{\begin{tikzpicture}[overlay]
		\node[draw=Horizon!60,line width=2.5pt,circle,fill=Horizon!25,font=\sffamily\bfseries,inner sep=4pt,outer sep=0pt] at (-19pt,5pt){\textcolor{Horizon}{Nota}};\end{tikzpicture}} % Blue Nota in a circle
\advance\baselineskip -1pt}{\end{list}\vskip5pt} % Tighter line spacing and white space after remark
\usepackage{graphicx}

\usepackage{titling}

\input{colores.tex}

%%Producir signos de cita textual``Asignatura''

\newcommand{\xmark}{\ding{55}} %%%comando que genera la tachita

\newenvironment{MyColorPar}[1]{%
\leavevmode\color{#1}\ignorespaces%
}{%
}%

%%%%%%%%%%%%%%%%%%% Título de la tarea, Nombre de alumno

\title{Clase \# 35 - Mecánica Teórica}
\author{\emph{JA} $\boldsymbol{\mid}$}
\date{\today}

%%%%%%%%%%%%%%%%%%% Datos de la Materia

\usepackage{fancyhdr}
\fancypagestyle{plain}{%  the preset of fancyhdr 
\fancyhf{} % clear all header and footer fields
\fancyfoot[R]{
\includegraphics[width=2cm]{LogoFCFMBUAP (1).png}}
\fancyfoot[L]{{\bfseries{\thedate{}}} a las {\bfseries{\mycurrenttime{}}} horas{} (GMT-6, H. Puebla de Zaragoza, Pue)}
\fancyhead[L]{Mecánica teórica}
\fancyhead[R]{\theauthor}
}
\makeatletter
\def\@maketitle{%
\newpage
\null
\vskip 1em%
\begin{center}%
\let \footnote \thanks
{\LARGE \@title \par}%
\vskip 1em%
%{\large \@date}%
\end{center}%
\par
\vskip 1em}
\makeatother

\usepackage{lipsum}  
\usepackage{cmbright}


%%%%%%%%%%%%%%%%%%%%%%%%%%%%%%%%%%%%%%%%%%%%%%%%%%%%%%%%%%%%%%%%%%%%%%%%%%%%%%%%%%%%%%%%%%%%%%%%%%%%%%%%%%%%%%%%%%%%%%%%%%%%%%%%%%%%%%%%%%%%%%%%%%%%%%%%%%%%%%%%%%%%%%%%%%%%%%%%%%%%%%%%%%%%%%%%%%%%%%

\begin{document}



\maketitle


%%%%%%%%%%%%%%%%%%% Datos del profesor, campus y aula

\begin{mybox}
\noindent\begin{tabular}{@{}ll}
	{\bfseries{Profesor}} & Dr. Gilberto Silva Ortigoza\\
	\textcolor{prussianblue}{\bfseries{Campus}} / \textcolor{trueblue}{\bfseries{Facultad}}  & \textcolor{prussianblue}{\bfseries{C.U. BUAP}} /  \textcolor{trueblue}{\bfseries{FCFM}} \\
	\textcolor{prussianblue}{\bfseries{Edificio}} /  \textcolor{trueblue}{\bfseries{EMA4}}    & 	\textcolor{prussianblue}{\bfseries{FM9}} /  \textcolor{trueblue}{\bfseries{302}}   
\end{tabular} 
\end{mybox} \vspace{1cm}

\tableofcontents

\vspace{2cm}

\section{Clase 35: Formulación Hamiltoniana}

\begin{figure}[H]
    \centering
	    \includegraphics[width=0.7\textwidth]{1.png}
    \caption{Representación esquemática de la construcción de la función Hamiltoniana a partir de la Lagrangiana \(L(q_i, \dot{q}_i, t)\). 
    Se definen los momentos generalizados \(p_i = \frac{\partial L}{\partial \dot{q}_i}\), 
    y se realiza la transformación de Legendre que conduce a 
    \(H(q_i, p_i, t) = \sum_{i=1}^{n} p_i \dot{q}_i - L(q_i, \dot{q}_i)\).}
    \label{fig:hamiltoniana_transformacion}
\end{figure}

Partimos de una Lagrangiana general dependiente de las coordenadas generalizadas, sus velocidades y el tiempo:
\[
L = L(q_i, \dot{q}_i, t)
\]

Mediante la \textbf{transformada de Legendre} definimos el momento canónico asociado a cada coordenada:
\[
p_i \equiv \frac{\partial L}{\partial \dot{q}_i}
\]

De modo que, si existen \(n\) coordenadas, tenemos el conjunto:
\[
p_1 = \frac{\partial L}{\partial \dot{q}_1}, \quad 
p_2 = \frac{\partial L}{\partial \dot{q}_2}, \quad 
\dots, \quad
p_n = \frac{\partial L}{\partial \dot{q}_n}
\]

Lo anterior permite expresar las velocidades como funciones de las variables \((q_i, p_i, t)\):
\[
\dot{q}_i = \dot{q}_i(q_i, p_i, t)
\]

Así, la \textbf{función Hamiltoniana} se define como:
\[
H(q_i, p_i, t) \equiv \sum_{i=1}^{n} p_i\,\dot{q}_i(q_i, p_i, t) - L\big(q_i, \dot{q}_i(q_i, p_i, t)\big)
\]

\noindent
De esta forma, el sistema de ecuaciones de movimiento hamiltonianas viene dado por:

\[
\boxed{
\dot{q}_{i} = \frac{\partial H}{\partial p_{i}},
\qquad
\dot{p}_{i} = -\,\frac{\partial H}{\partial q_{i}}
}
\]
\[
\text{(con } i = 1, 2, \dots, n)
\]

Si la Lagrangiana no depende explícitamente del tiempo, entonces:
\[
\frac{\partial H}{\partial t} = -\,\frac{\partial L}{\partial t} = 0
\]
lo cual implica que la energía total \(H\) es una constante del movimiento.

\subsection{Sobre la función generatriz}

En la formulación hamiltoniana, aparece una función \(S\) denominada \textbf{función generatriz}, que está relacionada con las transformaciones canónicas. 
Esta función puede emplearse para obtener un nuevo conjunto de variables \((Q_i, P_i)\) en el cual la forma de las ecuaciones de Hamilton se conserva.

En general, no podemos encontrar una transformación de punto que haga que la Lagrangiana se anule de manera idéntica; sin embargo, 
mediante la función generatriz \(S\) sí es posible transformar la descripción del sistema manteniendo la estructura del formalismo.

\subsection{Ejemplo introductorio}

Partimos de un problema en una dimensión, y lo extendemos al caso de dos dimensiones con un sistema libre. 
El vínculo entre ambas formulaciones —lagrangiana y hamiltoniana— se revela aquí con claridad: 
el paso de \(L\) a \(H\) no solo reescribe la dinámica, sino que revela las simetrías profundas del sistema.

\[
\text{¡Es realmente notable cómo la estructura se conserva bajo la transformación de Legendre!}
\]

\subsection{Sobre la singularidad de la Lagrangiana y la transformada de Legendre}

Sea la Lagrangiana de un sistema
\[
L = L(q_i, \dot{q}_i, t)
\]
y definamos la \textbf{matriz hessiana} (también llamada \emph{matriz de Legendre}) como
\[
W_{ij}(q,\dot q,t) \equiv \frac{\partial^2 L}{\partial \dot q_i\,\partial \dot q_j},
\qquad
\Delta \equiv \det(W_{ij})
\]

\paragraph{Caso no singular (regular).}
Si \(\Delta \neq 0\), la relación entre las velocidades y los momentos canónicos
\[
p_i = \frac{\partial L}{\partial \dot{q}_i}
\]
puede invertirse mediante el teorema de la función inversa, de modo que existen funciones
\[
\dot q_i = \dot q_i(q,p,t)
\]
que permiten definir sin ambigüedad la \textbf{transformada de Legendre}:
\[
H(q,p,t) = \sum_{i=1}^{n} p_i\,\dot q_i(q,p,t) - L\big(q,\dot q(q,p,t),t\big)
\]
En este caso, el formalismo hamiltoniano es completamente equivalente al lagrangiano:
\[
\boxed{
\dot q_i = \frac{\partial H}{\partial p_i},
\qquad
\dot p_i = -\frac{\partial H}{\partial q_i}
}
\]

\paragraph{Caso singular (degenerado).}
Si, en cambio, \(\Delta = 0\), la Lagrangiana se denomina \textbf{singular} o \textbf{degenerada}.
En este escenario, no todas las velocidades pueden resolverse en términos de \((q,p,t)\).
Surgen entonces relaciones entre las variables canónicas:
\[
\phi_a(q,p,t) = 0
\]
denominadas \textbf{restricciones primarias}.
Esto implica que el número de grados de libertad efectivos es menor que el número de coordenadas generalizadas \(n\).

El formalismo adecuado para tratar estos casos es el \textbf{formalismo de Dirac–Bergmann},
que clasifica las restricciones en:
\begin{itemize}
    \item \textbf{Primarias} y \textbf{secundarias}, según si provienen directamente de la definición de \(p_i\) o de la consistencia temporal;
    \item \textbf{De primera clase} o \textbf{de segunda clase}, según si generan simetrías (\emph{gauge}) o fijan valores concretos.
\end{itemize}

El número de grados de libertad físicos en el espacio de fase viene dado por:
\[
\boxed{ \text{gdl} = n - F - \tfrac{S}{2} }
\]
donde \(F\) es el número de restricciones de primera clase y \(S\) el de segunda clase.


Consideremos
\[
L = \tfrac{m}{2}\,(\dot q_1 - \dot q_2)^2
\]
Sus momentos canónicos son:
\[
p_1 = m(\dot q_1 - \dot q_2),
\qquad
p_2 = -m(\dot q_1 - \dot q_2)
\]
De donde se obtiene una restricción primaria:
\[
p_1 + p_2 = 0
\]
La matriz hessiana en este caso es:
\[
W_{ij} = m
\begin{pmatrix}
1 & -1\\[3pt]
-1 & 1
\end{pmatrix},
\qquad
\det W_{ij} = 0
\]
Por tanto, el sistema es singular y posee una restricción que reduce su número efectivo de grados de libertad.


Una Lagrangiana singular indica la presencia de \textbf{simetrías o dependencias entre las coordenadas},
por lo que las variables iniciales \((q_i, \dot q_i)\) no son independientes.
El paso a la formulación hamiltoniana, en estos casos, exige identificar dichas restricciones
y construir el \textbf{Hamiltoniano extendido} con multiplicadores de Lagrange:
\[
H_T = H + \lambda_a\,\phi_a(q,p,t)
\]
garantizando que la dinámica preserve las condiciones de consistencia en el tiempo.


La regularidad de la Lagrangiana —determinada por el valor de \(\det W_{ij}\)— 
es el criterio que decide si la \textbf{transformada de Legendre} puede aplicarse directamente o no.
Cuando la Lagrangiana es singular, el formalismo de Dirac provee las herramientas necesarias
para formular la dinámica reducida y determinar correctamente los grados de libertad físicos.


\subsection{¿Cuántos grados de libertad tiene el fotón?}

Si queremos estudiar la mecánica hamiltoniana, debemos estudiar formalismo de Dirac.

\subsection{Nos fijamos en ele ejemplo}


\[ L = \frac{m}{2} \left(  \dot{x}^{2} +  \dot{y}^{2} +  \dot{z}^{2} \right) \]

Tenemos:


\[ p_{x} = \frac{\partial L}{\partial \dot{x}} = m\dot{x} \quad \dot{x} = \frac{p_{x}}{m} \]

Y así con las demás componentes:

es decir:

\[  \boxed{ \dot{y} = \frac{ p_{y}}{m}} \quad \boxed{ \dot{z} = \frac{ p_{z}}{m}}   \]


El hamiltoniano no depende de las velocidades generalidades, no viven en el espacio de configuraciones, debemos escribirlas en posiciones iniciales, momentos iniciales, tiempos iniciales.


\[ H = \dot{x} p_{x} + \dot{y}p_{y} + \dot{z}p_{y} - \frac{1}{2}m\left( \dot{x}^{2} + \dot{y}^{2} + \dot{z}^{2}   \right)\]

\begin{figure}[H]
    \centering
	    \includegraphics[width=0.7\textwidth]{2.png}
    \caption{.}
    \label{fig:hamiltoniana_transformacion}
\end{figure}

\subsection{Calculando las ecuaciones de hamilton}

Observemos que $x,y,z$ son coordenadas cíclicas en la lagrangiana entonces en la hamiltoniana no van a aparecer.

\begin{figure}[H]
    \centering
	    \includegraphics[width=0.7\textwidth]{3.png}
    \caption{.}   
    \label{fig:hamiltoniana_transformacion}
\end{figure}

Estas tres ecuaciones ne forma vectorial dicen:

\[ \vec{r} = \vec{r}_{0} + \frac{\vec{P_{0}}}{m}t \]

momento:

\[ \vec{p}(t) = \vec{P_{0}}(t)  \]

\subsection{Ejemplo: Partícula de masa m y carga q sujeta la interacción dada por un campo magnético uniforme}

\begin{figure}[H]
    \centering
	    \includegraphics[width=0.7\textwidth]{4.png}
    \caption{.}   
    \label{fig:hamiltoniana_transformacion}
\end{figure}

El potencial vectorial $A$ no es único. 

Podría se un problema de examen.

Tenemos:

\[  p_{x} = \frac{\partial L}{\partial \dot{x}} = m\dot{x} - \frac{m\omega_{0}}{2}y \]

\[  p_{y} = \frac{\partial L}{\partial \dot{y}} = m\dot{y} + \frac{m\omega_{0}}{2}x \]

\[  p_{z} = \frac{\partial L}{\partial \dot{z}} = m\dot{z} \]


La idea es la siguiente, hacer los despejes correspondientes:


\[ \dot{x} = \frac{p_{x}}{m} + \frac{\omega_{0}}{2} y\]

\[\dot{y} = \frac{p_{y}}{m} - \frac{\omega_{0}}{2}x  \]

\[ \dot{z} = \frac{p_{z}}{m} \]


NO siempre vamos a tener problemas en donde se puedan despejar de esta forma las velocidades. puede ser que en el examen tengamos que calcularlos de diferente manera.

Tenemos así la \textbf{Hamiltoniana}:


\[ H = \dot{x}p_{x}+\dot{y}p_{y}+\dot{z}p_{z}-L \]

\[H = p_{x} \left(   \right)  \]

\begin{figure}[H]
    \centering
	    \includegraphics[width=0.8\textwidth]{5.png}
    \caption{.}   
 \label{fig:hamiltoniana_transformacion}
\end{figure}

Si realizamos las cuentas llegamos al \textbf{hamiltoniano} siguiente:


\[  H(\vec{r},\vec{p}, t) = \frac{1}{2m} \left( p^{2}_{x} + p^{2}_{y} + p^{2}_{z} \right) \ + \frac{\omega_{0}}{2} \left( yp_{x}-xp_{y}  \right) +\frac{m\omega_{0}^{2}}{8} \left( x^{2} + y^{2} \right) \]

Calculamos las ecuaciones canónicas de hamilton:

Considerando que tenemos \[ q_{1} = x; \quad  q_{2} = y; \quad q_{3} = z\]


La componente $p_{x}$,$p_{y}$,$p_{z}$ son generadores de traslaciones

\[ \dot{q}_{i} = \frac{\partial H}{\partial p_{i}}; \quad \dot{p}_{i} =- \frac{\partial H}{\partial q_{i}} \]

\begin{figure}[H]
    \centering
	    \includegraphics[width=0.7\textwidth]{6.png}
    \caption{.}   
 \label{fig:hamiltoniana_transformacion}
\end{figure}

Este es un sistema de 6 incognitas y es lo que tenemos que resolver:

\section{Principio de Hamilton}

\begin{figure}[H]
    \centering
	    \includegraphics[width=0.7\textwidth]{7.png}
    \caption{.}   
 \label{fig:hamiltoniana_transformacion}
\end{figure}


Si las contricciones son holonomas, la trayectoria que corresponde a la diámica del sistema mećanico es aquellla que hace estarionaria a la acción denretro de todo el conjunto de trayentoria, si es que hacmeos variaciones infinitesimales, es decir difieren muy poco una de otra las trayectorias.

\[ \delta{S}^{1} =0 \quad \Longleftrightarrow \quad \frac{d}{dt} \left(  \frac{\partial L}{\partial \dot{q}_{i}} \right) - \frac{\partial L}{\partial q_{i}} =0 \]

\subsection{Espacio fase}

\begin{figure}[H]
    \centering
	    \includegraphics[width=0.7\textwidth]{8.png}
    \caption{.}   
 \label{fig:hamiltoniana_transformacion}
\end{figure}

¿Cómo determinamos la curva? 


Se puede determinar de la siguiente manera:


Existe otra forma de escribir la lagrangiana:

\subsubsection{Principio de Hamilton modificado}

\begin{figure}[H]
    \centering
	    \includegraphics[width=0.7\textwidth]{9.png}
    \caption{.}   
 \label{fig:hamiltoniana_transformacion}
\end{figure}

Tenemos:

\[ L = L (q_{i},\dot{q}_{i}, t ) \]


Tenemos la nueva Lagrangiana:


\begin{figure}[H]
    \centering
	    \includegraphics[width=0.7\textwidth]{10.png}
    \caption{.}   
 \label{fig:hamiltoniana_transformacion}
\end{figure}

En la acción básicas depende de cordenada generalizada

\begin{figure}[H]
    \centering
	    \includegraphics[width=0.7\textwidth]{11.png}
    \caption{Ecuaciones de Hamilton}   
 \label{fig:hamiltoniana_transformacion}
\end{figure}

Tenemos haciendo el análisis para $L_{H}$:

\begin{figure}[H]
    \centering
	    \includegraphics[width=0.7\textwidth]{12.png}
    \caption{Ecuaciones de Hamilton}   
 \label{fig:hamiltoniana_transformacion}
\end{figure}

Tenemos:


\[ L_{H} = \dot{q}p - H(q,p,t)  \]


\[L^{\prime} = L + \frac{dF}{dt} (q,t)  \]

\[ L^{\prime}_{H} = L_{H} + \frac{dF}{dt}(q,p,t) \]

\[  = L_{H} + \frac{\partial F}{\partial q}\dot{q}+\frac{\partial F}{\partial p}\dot{p} + \frac{\partial F}{\partial t} \]


Formulación hamiltoniana tiene su desarrollo por sí mismo, en principio es independiente de la formulación hamiltoniana.

\end{document}